% chapters/05/sub-cpmk.tex
\begin{subcpmkbox}
Bab ini menuntaskan \SUBCPMK{1.2} secara utuh. Ketiga indikatornya diperiksa,
dan buktinya berupa program yang menggerakkan objek dengan kecepatan tetap
terlepas dari laju bingkainya.

\vspace{0.5em}
\textbf{\SUBCPMK{1.2} Merancang dan mengimplementasikan gerak objek berbasis
\istilah{delta time} dan vektor pada Pygame} --- ranah C3, \CPMK{1},
\CPL{03}.

\vspace{0.5em}
{\setlength{\tabcolsep}{4pt}\footnotesize%
\begin{tabularx}{\textwidth}{@{}C{1.9cm} L{4.0cm} L{2.5cm} Y@{}}
\toprule
\headrow \textbf{Indikator} & \textbf{Rumusan} & \textbf{Dibahas pada} & \textbf{Bukti} \\
\midrule
\ind{a} & Menulis kode gerak menggunakan \code{pygame.Vector2} dan
\istilah{delta time} &
Bagian~\ref{sec:posisi-vektor} sampai Bagian~\ref{sec:siklus-gerak} &
Program \code{ch05_gerak.py} pada repositori, beserta penjelasan baris
perhitungan perpindahannya. \\
\addlinespace
\ind{b} & Menunjukkan kecepatan objek tetap konsisten pada 30 FPS dan 60 FPS &
Bagian~\ref{sec:gerak-selang-waktu} &
Tabel pengukuran pada lembar kerja praktikum: jarak tempuh dan kecepatan
terukur pada tiga laju bingkai. \\
\addlinespace
\ind{c} & Menyinkronkan posisi bertipe \istilah{float} dengan \code{Rect}
bertipe bilangan bulat &
Bagian~\ref{sec:selaras-rect} &
Potongan kode penyelarasan posisi beserta penjelasan mengapa posisi tidak boleh
disimpan sebagai bilangan bulat. \\
\bottomrule
\end{tabularx}}
\end{subcpmkbox}
